\documentclass[conference]{IEEEtran}
\IEEEoverridecommandlockouts
\usepackage{cite}
\usepackage{amsmath,amssymb,amsfonts}
\usepackage{graphicx}
\usepackage{textcomp}
\usepackage{xcolor}
\usepackage{booktabs}
\usepackage{url}

% Registered results resolve through docs/research/result-traceability.md
% and docs/research/findings.md. The preliminary encoder comparison uses
% runs/encoders-20261022/checkpoints/stream-lipsync-avhubert-history.json.

\begin{document}

\title{Multimodal Audio-Visual Deepfake Detection:\\
Validity Checks Catch a Detector That Looks Right and Is Not}

\author{\IEEEauthorblockN{Vaibhav Shiroorkar}
\IEEEauthorblockA{\textit{Dept. of Electronics and Computer Engineering} \\
\textit{KJ Somaiya School of Engineering}\\
Mumbai, India \\
v.shiroorkar@somaiya.edu}
\and
\IEEEauthorblockN{Adithya Jithesh}
\IEEEauthorblockA{\textit{Dept. of Electronics and Computer Engineering} \\
\textit{KJ Somaiya School of Engineering}\\
Mumbai, India \\
adithya.j@somaiya.edu}
\and
\IEEEauthorblockN{Shivansh Gupta}
\IEEEauthorblockA{\textit{Dept. of Electronics and Computer Engineering} \\
\textit{KJ Somaiya School of Engineering}\\
Mumbai, India \\
shivansh.gupta@somaiya.edu}
\and
\IEEEauthorblockN{Tejashree Phatak}
\IEEEauthorblockA{\textit{Dept. of Electronics and Computer Engineering} \\
\textit{KJ Somaiya School of Engineering}\\
Mumbai, India \\
tejashree@somaiya.edu}
}

\maketitle

\begin{abstract}
Lip-sync forgeries change only the mouth and the voice, so a detector that
reads pixels alone sees mostly genuine video. Combining audio and video is the
natural answer. We built that system: visual streams, a voice stream, and
cross-attention streams that compare the mouth with the sound and the face
with the voice, fused at the feature level. Then we tested whether it does
what its parts are named for. We ask whether cue-specific fusion generalizes
better than a strong visual baseline under source-disjoint, shortcut-controlled
evaluation. The historical fusion comparison requires a nested rerun. Historical fusion comparisons found an in-domain gain and no cross-corpus
gain, but their outer-holdout labels selected branch checkpoints. Those
comparisons require nested cross-fitting before they support a conclusion. The stronger result is how close the project came to reporting
the opposite. A cross-attention stream reached 0.9991 ROC-AUC while its
attention stayed at the uniform value, so its accuracy is real and its named
mechanism is absent. The served visual models partly read the video codec:
re-encoding genuine clips with the codec of the largest fake category makes
93.3 percent of them look fake. Compression augmentation with codec-swapped
checkpoint selection cost 0.19 of cross-corpus ROC-AUC.
A checkpoint chosen on validation loss scored below
chance, and the training seed alone moves cross-corpus ROC-AUC by 0.1078.
Each was caught by a quantity recorded beside the training objective and never
optimized against it. We release the evaluation protocol, the measurements and
their failures.
\end{abstract}

\begin{IEEEkeywords}
Audio-visual deepfakes, cross-attention, cross-corpus generalization,
feature-level fusion, media forensics, shortcut learning.
\end{IEEEkeywords}

\paragraph{Evidence validity}
The fusion, LAV-DF and PE/CLIP tables retain historical measurements.
The old OOF protocol used outer-holdout labels for epoch selection.
Explicit zero-start LAV-DF windows could overlap manipulated spans.
PE/CLIP features used incorrect normalization.
These results require reruns and do not establish corrected performance.

\section{Introduction}

Early deepfakes swapped whole faces. The most convincing forgeries today are
narrower. An attacker generates speech with text-to-speech or voice
conversion, then regenerates only the mouth to match it. The background, the
lighting and the upper face stay genuine. A detector that reads pixels alone
sees a video that is mostly real.

The natural response is to read both modalities and look for disagreement: a
mouth shaping one sound while the audio carries another, or a voice whose
affect the face does not show. We built a system on that idea. Visual streams
read sequences of face crops. A voice stream reads the audio. Two
cross-attention streams compare the mouth with the sound and the face with the
voice. A learned head fuses all of them at the feature level.

The research question was fixed before the experiments and anticipated a
negative answer: under source-disjoint and shortcut-controlled evaluation,
does cue-specific fusion of visual artifacts, audio spoofing and mouth-audio
alignment generalize better than a strong visual baseline? The historical
comparison suggested no gain, but requires a nested cross-fitting rerun.

The more useful finding is what we had to check to say so. Four results would
each have passed unnoticed under an evaluation built on accuracy tables alone:

\begin{enumerate}
\item A cross-attention stream reached near-perfect accuracy with attention
that never left its uniform starting value.
\item The visual models partly learned the video codec, because one dataset
encodes its largest fake category differently from its genuine clips.
\item A selection rule kept a checkpoint that scored below chance.
\item The training seed moved the cross-corpus result more than any
architectural choice we tested.
\end{enumerate}

Each was caught by a quantity recorded next to the objective and never
optimized against it. That practice, rather than the architecture, is what we
contribute. A second track applies it to video generated whole, where a frozen
encoder with a linear probe is scored on generators it never saw.

\section{Related Work}

Audio-visual deepfake corpora such as FakeAVCeleb~\cite{fakeavceleb} and
LAV-DF~\cite{lavdf} pair manipulated faces with manipulated voices, and
cross-dataset evaluation on Celeb-DF~\cite{celebdf} and DFDC~\cite{dfdc} after
training on FaceForensics++~\cite{ffpp} is the standard test of transfer.
Mouth-audio synchronization has been modelled since SyncNet~\cite{syncnet},
and AV-HuBERT~\cite{avhubert} learns a joint audio-visual speech
representation from unlabelled video. Self-blended images~\cite{sbi} train a
face detector on blending artifacts rather than on any one generator, and a
frozen vision-language encoder with a linear probe transfers well across image
generators~\cite{univfd}. Shortcut learning~\cite{shortcut}, where a model
solves the benchmark by a cue the task did not intend, is the failure most of
our checks look for. A concurrent audit of deepfake benchmarks with frozen
self-supervised features~\cite{audit} covers ground close to our
region-of-interest result.

We do not claim multimodal fusion is new. We claim the evidence practices
around it are underspecified in a way that lets a stream with an inert
mechanism pass as a working one, and we give instances.

\section{Method}

\subsection{Two systems, one evaluation path}

Design A has three branches. EfficientNet-B0~\cite{efficientnet} with a GRU
reads 16 face crops of $224\times224$ pixels, detected and tracked with
MTCNN~\cite{mtcnn}. Wav2Vec2~\cite{wav2vec2} with attentive pooling reads 4~s
of 16~kHz audio for spoofing. A framewise ResNet-18 mouth encoder with a
second Wav2Vec2 reads a 2~s window of 50 mouth crops for mouth-audio
alignment. Fusion is regularized logistic regression over Platt-calibrated
branch logits and three quality features, fitted only on out-of-fold rows.

Design B replaces the branches with configurable streams and the late head
with a feature-level head. A visual stream is a backbone (EfficientNet-B0, or
a frozen DINOv3 ViT-S/16~\cite{dinov3}) followed by a temporal model. An
audio-visual stream is cross-attention between temporal tokens:
\begin{equation}
A = \operatorname{softmax}\!\left(\frac{(X_a W_Q)(X_v W_K)^\top}{\sqrt{d}}\right),
\qquad \tilde{X} = A\,(X_v W_V),
\end{equation}
where the audio tokens $X_a$ ask and the video tokens $X_v$ answer. The stream
embeds the residual $X_a - \tilde{X}$, the part of the sound the mouth did not
account for, averaged over time. The lip-sync stream reads the mouth window and
its audio. The emotion stream reads the face crops and the voice. The fusion
head projects each stream's 256-dimensional embedding to a common width, masks
absent streams after the projection so that an absent stream contributes
exactly zero, and classifies the concatenation with a multi-layer perceptron.

Both systems are scored through the same code. The largest gap in our results
was first suspected to be an evaluation artifact, and scoring the Design A
checkpoint through the Design B path reproduced its number to four decimals.

\subsection{A validity check recorded and never optimized}

A stream that aligns mouth motion with sound should concentrate its attention
near the diagonal of $A$, since sound arrives at a near-fixed offset from the
articulation that produced it. We record the diagonal mass
\begin{equation}
m = \frac{1}{T_q}\sum_{i,j} A_{ij}\,\mathbb{1}\!\left[\,\left|\tfrac{T_k}{T_q}\,i - j\right| \le 1\right]
\end{equation}
every epoch and never put it in the loss. A measure that is optimized stops
being evidence about the mechanism and becomes another thing the model fits.
Its chance value depends on the query length: uniform attention gives 0.0592
at 50 steps and 0.3438 at 8.

\subsection{Abstention and provenance}

A clip with no stable face track yields no visual view. The pipeline reports
it as unavailable rather than substituting a full frame or dropping it, and
coverage is reported beside every metric. Splits are frozen by source identity.
The view cache is addressed by a hash of the preprocessing settings and code
version, and a model is refused on a cache built differently. Registered
results carry artifact SHA-256 hashes. The preliminary encoder comparison
below uses a training history and has no held-out evaluation yet.

\section{Experimental Setup}

Track one trains on FakeAVCeleb~\cite{fakeavceleb}, split 70/15/15 by source
identity. The cross-corpus test is DFDC~\cite{dfdc}, the only audio-bearing
corpus here not built on VoxCeleb2. Celeb-DF-v2~\cite{celebdf} and the
evaluation-only MNW benchmark~\cite{mnw} are external checks. LAV-DF~\cite{lavdf}
is used to test the cross-modal streams under a stricter construction.
FaceForensics++~\cite{ffpp} c23 serves as a second training corpus, following
the published cross-dataset protocol.

Track two uses DF26~\cite{df26}: 271 real clips and 2,420 from seven
generators, four commercial and three open source. A frozen encoder with a
logistic probe trains on six generators and is scored on the seventh.

The protocol fixes training seeds 17, 29 and 43. All Design A branches and the
fusion head have all three. Every Design B stream has seed 17 alone, and seed
17 is the most favourable of the three on DFDC, so a single-seed cross-corpus
figure here is the best draw rather than a central one.

Confidence intervals are 95 percent bootstrap intervals~\cite{bootstrap}
clustered on source identity, because a reader's question is what happens on a
different set of speakers. A resample that lands on one class is skipped and
counted, never scored 0.5. Comparisons between two systems are paired: the
same identities are resampled for both and the difference is taken within each
resample. DF26 intervals resample clips or source videos, since generated
clips have no speaker.

\section{Results}

\subsection{The cross-modal streams do not read correspondence}

The emotion stream reads 0.9991 ROC-AUC on held-out FakeAVCeleb clips, the best
of the five streams. Table~\ref{tab:diag} shows its diagonal mass and the
lip-sync stream's across every epoch of three training runs.

\begin{table}[htbp]
\caption{Diagonal attention mass against its uniform value}
\label{tab:diag}
\centering
\begin{tabular}{lrrl}
\toprule
Stream & Steps & Uniform & Measured \\
\midrule
emotion & 8 & 0.3438 & 0.3434 to 0.3457, 9 epochs \\
lip-sync & 50 & 0.0592 & 0.0592 to 0.0595, 4 epochs \\
lip-sync, lower encoder LR & 50 & 0.0592 & 0.0592 to 0.0595, 6 epochs \\
\bottomrule
\end{tabular}
\end{table}

Emotion's ROC-AUC rose from 0.9847 to 0.9991 over nine epochs while its
attention did not move. The accuracy is real and the mechanism the stream is
named for is not operating.

FakeAVCeleb's fakes are whole-clip manipulations, so a model can separate them
by recognising the generator without comparing mouth and sound. LAV-DF removes
that route: we cut two windows from the same file, one containing the
manipulated span and one avoiding it, so speaker, lighting, microphone and
codec are held fixed. A lip-sync stream trained on those pairs reads 0.9969
ROC-AUC on 1,161 held-out windows with diagonal mass 0.0592 against a uniform
0.0600. The matched windows control file-level cues, but local manipulation
artifacts can still separate the classes. Near-uniform diagonal mass does not
establish that the model learned temporal correspondence.

\subsection{The visual models partly read the video codec}

FakeAVCeleb's genuine videos are H.264 at about 150~kbps. Every face-swap plus
Wav2Lip fake, its largest fake category, is MPEG-4 Part~2 at about 330~kbps.
Re-encoding 75 test clips per row and scoring them with the served ensemble at
threshold 0.5 gives Table~\ref{tab:codec}.

\begin{table}[htbp]
\caption{Called fake after re-encoding (75 clips per row)}
\label{tab:codec}
\centering
\begin{tabular}{llr}
\toprule
Clips & Codec & Called fake \\
\midrule
genuine & original H.264 & 0.0\% \\
genuine & re-encoded H.264 (control) & 0.0\% \\
genuine & re-encoded MPEG-4 & \textbf{93.3\%} \\
swap plus Wav2Lip & original MPEG-4 & 100.0\% \\
swap plus Wav2Lip & re-encoded H.264 & \textbf{30.7\%} \\
face swap, H.264 & re-encoded MPEG-4 & 100.0\% \\
\bottomrule
\end{tabular}
\end{table}

The control decides the reading. Re-encoding a genuine clip with its own codec
changes its pixels and moves nothing. Changing only the codec moves the verdict
in both directions. Decoding recovers the pixels the encoder wrote,
compression artifacts included, so a model can learn the encoder.

\subsection{Training against the codec cue costs transfer}

We trained four EfficientNet-B0 recipes at seed 17, each changing one thing
against the one before: the shipped recipe selected on ROC-AUC; added random
JPEG, H.264 and MPEG-4 round trips on both classes with checkpoint selection on
codec-swapped validation clips; added self-blended fakes~\cite{sbi}; and added
the FaceForensics++ training split.

\begin{table}[htbp]
\caption{Visual recipes against the codec cue (ROC-AUC)}
\label{tab:recipes}
\centering
\begin{tabular}{lrrr}
\toprule
Recipe & FakeAVCeleb & DFDC & Celeb-DF \\
\midrule
shipped & 0.9988 & \textbf{0.7611} & 0.6682 \\
+ compression & 0.9959 & 0.5720 & 0.6525 \\
+ self-blending & 0.9978 & 0.6496 & 0.6762 \\
+ FaceForensics++ & 0.9904 & 0.6756 & 0.6509 \\
\bottomrule
\end{tabular}
\end{table}

Every modified recipe scored below the shipped recipe on DFDC
(Table~\ref{tab:recipes}), compression most, by
0.1891 with the intervals apart. Self-blending and the larger corpus won back
part of it and still sit 0.0855 below the shipped recipe. Celeb-DF separates
none of the four by interval overlap alone. Compression augmentation and
checkpoint selection changed together, so this comparison does not isolate
the codec cue's contribution to transfer. The shipped recipe, retrained,
reproduced the original checkpoint exactly, epoch for epoch.

\subsection{Fusion does not beat the visual branch across corpora}

Table~\ref{tab:fusion} gives the paired difference between the fused system and
a head fitted on the visual branch plus the same quality features, each seed on
its own cross-fitted fold models.

\begin{table}[htbp]
\caption{Fusion minus visual, paired, 95\% interval}
\label{tab:fusion}
\centering
\begin{tabular}{lll}
\toprule
Seed & FakeAVCeleb test & DFDC \\
\midrule
17 & $+0.0159$ [$+0.0129$, $+0.0190$] & $-0.0085$ [$-0.0196$, $+0.0027$] \\
29 & $+0.0170$ [$+0.0143$, $+0.0195$] & $-0.0215$ [$-0.0371$, $-0.0052$] \\
43 & $+0.0184$ [$+0.0159$, $+0.0209$] & $-0.0417$ [$-0.0698$, $-0.0129$] \\
\bottomrule
\end{tabular}
\end{table}

The in-domain gain holds at every seed, and it comes almost entirely from the
75 clips whose video is genuine and whose voice was replaced, which no visual
branch can detect. Across corpora fusion is never better and is measurably
worse at two seeds. In Design B, every subset of the five streams was fitted
on the same rows: the best DFDC combination, emotion plus lip-sync, reads
0.6076, and fusing all five reads 0.5411, below three of its own inputs. The
audio branch is at chance across corpora, 0.5054 [0.4747, 0.5371], and the
sync branch's anomaly score is below 0.5 at all three seeds. One branch
transfers and two do not.

\subsection{Seed spread exceeds the measured fusion differences}

\begin{table}[htbp]
\caption{DFDC ROC-AUC by training seed}
\label{tab:seeds}
\centering
\begin{tabular}{lrrrr}
\toprule
System & 17 & 29 & 43 & Spread \\
\midrule
visual branch & 0.7611 & 0.6963 & 0.6878 & 0.0733 \\
fused system & 0.7500 & 0.6698 & 0.6422 & \textbf{0.1078} \\
\bottomrule
\end{tabular}
\end{table}

In-domain the seeds agree to three or four decimal places. Across corpora the
fused system ranges over 0.1078 (Table~\ref{tab:seeds}), thirteen times the
seed-17 fusion difference. Fusing three branches compounded the seed noise
rather than averaging it. Averaging the six final visual checkpoints
(EfficientNet-B0 and ConvNeXt-Tiny~\cite{convnext} at three seeds) with equal
weights reads 0.7739 [0.7401, 0.8064] on DFDC, $+0.0469$ [$+0.0286$, $+0.0666$]
over the EfficientNet seed mean and level with its best seed, and level with
every member on Celeb-DF. The ensemble grouping was chosen after comparing
three groupings on DFDC. Its DFDC gain therefore needs confirmation on an
untouched test partition.

\subsection{Choices that looked right on validation}

\textit{Selection on loss.} One visual stream reached its lowest validation
loss at epoch 1 with its backbone still frozen. That checkpoint scored 0.4668
ROC-AUC in-domain, below chance. Binary cross-entropy punishes a confident
mistake, so a model that ranks well while overconfident loses to one that
hedges and ranks badly. Every trainer now selects on validation ROC-AUC.

\textit{Frozen BatchNorm.} Freezing BatchNorm statistics during fine-tuning
raised validation from 0.5154 to 0.9058 and was first recorded as a clear win.
Live statistics cost 0.5430 of in-domain ROC-AUC relative to frozen statistics
and gained 0.1276 on DFDC. The arms ran different numbers of epochs.
With the budget held
fixed the DFDC difference is 0.0177, one seventh of the uncontrolled figure.
Live statistics act as an accidental domain adapter~\cite{adabn}, and the
in-domain cost is what survives the control.

\subsection{The operating point fails before the ranking does}

\begin{table}[htbp]
\caption{Visual branch under corruption of the cached view}
\label{tab:corruption}
\centering
\begin{tabular}{lrrr}
\toprule
Condition & ROC-AUC & Genuine called fake & Fakes missed \\
\midrule
clean & 0.9988 & 0 of 150 & 3 of 1,498 \\
JPEG 75 & 0.9934 & \textbf{76 of 150} & 1 \\
JPEG 50 & 0.9187 & \textbf{150 of 150} & 0 \\
\bottomrule
\end{tabular}
\end{table}

JPEG at quality 75 costs 0.0054 of ROC-AUC, which a table reads as robustness,
while taking false alarms from 0 to 51 percent (Table~\ref{tab:corruption}).
Audio noise at 20~dB SNR fails the other way: 0.0227 of ROC-AUC, and 857 of 859
fakes missed. Across corpora the model is also more confident when wrong than
when right: on DFDC the visual branch is wrong on 1,942 of 2,410 clips at 0.5,
with mean confidence 0.9244 when wrong against 0.8517 when right. No threshold
on that confidence drops errors faster than correct answers. At 0.5 the visual
branch calls 141 of 155 genuine Celeb-DF videos fake.

A Gaussian over frozen DINOv3 features of genuine training clips, scored by
Mahalanobis distance, never sees a fake and asks only whether the footage is
familiar. At the served setting it abstains on 8.7 percent of held-out
FakeAVCeleb genuine clips and intercepts 74.4 percent of genuine Celeb-DF
clips, between 49.4 and 83.1 percent over twenty reference splits. False
accusations on unfamiliar footage fall from about 91 percent to between a
quarter and a third.

\subsection{Open-world generated video}

Four shortcuts in DF26 were caught before any number was reported: scene
content, real clips shared across splits, sibling modes of one generator, and
audio presence, which alone predicts the label because every commercial clip
carries a track and almost no other clip does. Real clips were also the only
ones at 25 or 29.97 frames per second, so every clip was re-encoded to
$640\times360$ at 24 frames per second.

\begin{table}[htbp]
\caption{Unseen-generator ROC-AUC on DF26, frozen encoder and linear probe}
\label{tab:df26}
\centering
\begin{tabular}{lrl}
\toprule
Encoder & Macro & Paired against DINOv3 \\
\midrule
Perception Encoder~\cite{pe} & \textbf{0.7527} & $+0.0741$ [$+0.0121$, $+0.1317$] \\
CLIP ViT-B/16~\cite{clip} & 0.7454 & $+0.0668$ [$+0.0049$, $+0.1289$] \\
DINOv3 ViT-S/16 & 0.6786 & reference \\
\bottomrule
\end{tabular}
\end{table}

Perception Encoder leads on all seven held-out generators
(Table~\ref{tab:df26}), from 0.8496 on Veo~3.1 to 0.6542 on LTX~2.3. These
encoders were fed ImageNet-normalized pixels, which is right for DINOv3 and
wrong for the other two. These are measurements under mismatched preprocessing,
not lower bounds. The comparison remains provisional until features are
re-extracted with each encoder's own statistics.

The region of interest decides which design works. The same full-frame probe
on FaceForensics++, one manipulation family held out per arm, is at chance,
macro 0.5387. Face crops through the fine-tuned visual stream, retrained with
one family held out each time, lose 0.1231 macro against a model that saw all
six, and the graphics-based FaceSwap falls to 0.5195. Full frames suit video
generated whole; face crops suit a manipulated real face.

\subsection{Encoders the cross-modal design named}

The stream design named AV-HuBERT~\cite{avhubert} for lip-sync and
HSEmotion~\cite{hsemotion} for the face side of emotion. Neither runs in this
environment as released: AV-HuBERT needs fairseq, which pins an older PyTorch,
and HSEmotion ships a pickled model that current timm cannot execute. We
restated AV-HuBERT Base in plain PyTorch and loaded HSEmotion's weights into
the same EfficientNet-B0 architecture, and checked both against the official
code on identical inputs: features and logits agree to $2\times10^{-6}$.
AV-HuBERT reads once on the lips alone and once on the audio alone,
since its pretraining placed both in one space.

Aligned audio agrees with the lips better than audio shifted by half a second
on 98 percent of 128 genuine training clips. A lip-sync stream over these
features reaches 0.9416 validation ROC-AUC, against 0.7593 for the Wav2Vec2
stream. The selected checkpoint is epoch 2 of five, before encoder unfreezing
at epoch 3. Diagonal mass ranges from 0.05875 to 0.05919, close to the uniform
value of 0.0592. The validation gain provides no evidence of diagonal
alignment in the trained attention head. The shifted-audio check concerns the
pretrained features; it does not validate the downstream detector's mechanism.

This comparison is preliminary and uses seed 17. Held-out FakeAVCeleb, DFDC,
and fusion scores for the replacement encoders are pending. HSEmotion training
has no completed result in this comparison. Validation accuracy alone cannot
establish that either replacement improves cross-corpus detection.

\section{Limitations}

Every Design B stream has one training seed where the protocol asks for three,
and Table~\ref{tab:seeds} shows what a seed is worth across corpora. Design A
and B train on FakeAVCeleb, so their numbers compare against each other rather
than against published tables; the FaceForensics++ result follows the
published protocol at one seed. No training data comes from the wild. The
cross-modal streams do not measure correspondence, so the mechanism the
architecture is named for is unverified. One exit criterion of the research
design, that the sync offset interval exclude chance, was never measured,
because the branch's history recorded only a combined loss. Track one is close
to blind on fully generated video: 0 of 10 MNW clips from one talking-head
generator are detected. Track two rests on one corpus and one split, and DF26
cannot support audio work.

\section{Conclusion}

Cue-specific fusion did not beat the visual baseline across corpora on our
data, and fusing every stream was worse than three of its parts. That answers
the question as asked. The more useful result is how close the project came to
reporting the opposite: a stream at 0.9991 with an inert mechanism, a visual
model reading the codec, a regularizer that bought validation accuracy by
removing a domain adapter, and a selection rule that preferred a below-chance
checkpoint. Each was caught by a quantity recorded beside the objective and
never optimized against it, and the same practice found the project's own
errors, including a reported effect that shrank sevenfold once its control
ran. Reporting one cross-corpus number to four decimal places invites
comparisons the measurement cannot support when the seed alone moves it by
0.1078. A validity protocol needs mechanism checks and seed spread together.

\begin{thebibliography}{00}
\bibitem{fakeavceleb} H. Khalid, S. Tariq, M. Kim, and S. S. Woo,
``FakeAVCeleb: A novel audio-video multimodal deepfake dataset,'' in
\textit{Proc. NeurIPS Datasets and Benchmarks Track}, 2021.
\bibitem{lavdf} Z. Cai, K. Stefanov, A. Dhall, and M. Hayat, ``Do you really
mean that? Content driven audio-visual deepfake dataset and multimodal method
for temporal forgery localization,'' in \textit{Proc. DICTA}, 2022.
\bibitem{celebdf} Y. Li, X. Yang, P. Sun, H. Qi, and S. Lyu, ``Celeb-DF: A
large-scale challenging dataset for DeepFake forensics,'' in \textit{Proc.
CVPR}, 2020.
\bibitem{dfdc} B. Dolhansky, J. Bitton, B. Pflaum, J. Lu, R. Howes, M. Wang,
and C. Canton Ferrer, ``The DeepFake Detection Challenge (DFDC) dataset,''
arXiv:2006.07397, 2020.
\bibitem{ffpp} A. R\"ossler, D. Cozzolino, L. Verdoliva, C. Riess, J. Thies,
and M. Nie{\ss}ner, ``FaceForensics++: Learning to detect manipulated facial
images,'' in \textit{Proc. ICCV}, 2019.
\bibitem{syncnet} J. S. Chung and A. Zisserman, ``Out of time: Automated lip
sync in the wild,'' in \textit{Proc. ACCV Workshops}, 2016.
\bibitem{avhubert} B. Shi, W.-N. Hsu, K. Lakhotia, and A. Mohamed, ``Learning
audio-visual speech representation by masked multimodal cluster prediction,''
in \textit{Proc. ICLR}, 2022.
\bibitem{sbi} K. Shiohara and T. Yamasaki, ``Detecting deepfakes with
self-blended images,'' in \textit{Proc. CVPR}, 2022.
\bibitem{univfd} U. Ojha, Y. Li, and Y. J. Lee, ``Towards universal fake image
detectors that generalize across generative models,'' in \textit{Proc. CVPR},
2023.
\bibitem{shortcut} R. Geirhos, J.-H. Jacobsen, C. Michaelis, R. Zemel,
W. Brendel, M. Bethge, and F. A. Wichmann, ``Shortcut learning in deep neural
networks,'' \textit{Nature Machine Intelligence}, vol. 2, pp. 665-673, 2020.
\bibitem{audit} ``What do deepfake benchmarks measure? An audit using frozen
self-supervised representations,'' arXiv:2606.26384, 2026.
\bibitem{efficientnet} M. Tan and Q. V. Le, ``EfficientNet: Rethinking model
scaling for convolutional neural networks,'' in \textit{Proc. ICML}, 2019.
\bibitem{mtcnn} K. Zhang, Z. Zhang, Z. Li, and Y. Qiao, ``Joint face detection
and alignment using multitask cascaded convolutional networks,'' \textit{IEEE
Signal Process. Lett.}, vol. 23, no. 10, pp. 1499-1503, 2016.
\bibitem{wav2vec2} A. Baevski, Y. Zhou, A. Mohamed, and M. Auli, ``wav2vec 2.0:
A framework for self-supervised learning of speech representations,'' in
\textit{Proc. NeurIPS}, 2020.
\bibitem{dinov3} O. Sim\'eoni \textit{et al.}, ``DINOv3,'' arXiv:2508.10104,
2025.
\bibitem{mnw} Microsoft, Northwestern University, and WITNESS, ``MNW deepfake
detection benchmark,'' \url{https://github.com/microsoft/MNW}.
\bibitem{df26} ``DF26,'' dataset,
\url{https://huggingface.co/datasets/DF26/DF26}, 2026.
\bibitem{bootstrap} B. Efron and R. J. Tibshirani, \textit{An Introduction to
the Bootstrap}. New York: Chapman \& Hall, 1993.
\bibitem{convnext} Z. Liu, H. Mao, C.-Y. Wu, C. Feichtenhofer, T. Darrell, and
S. Xie, ``A ConvNet for the 2020s,'' in \textit{Proc. CVPR}, 2022.
\bibitem{adabn} Y. Li, N. Wang, J. Shi, X. Hou, and J. Liu, ``Adaptive batch
normalization for practical domain adaptation,'' \textit{Pattern Recognition},
vol. 80, pp. 109-117, 2018.
\bibitem{pe} D. Bolya \textit{et al.}, ``Perception Encoder: The best visual
embeddings are not at the output of the network,'' arXiv:2504.13181, 2025.
\bibitem{clip} A. Radford \textit{et al.}, ``Learning transferable visual
models from natural language supervision,'' in \textit{Proc. ICML}, 2021.
\bibitem{hsemotion} A. V. Savchenko, ``Facial expression and attributes
recognition based on multi-task learning of lightweight neural networks,'' in
\textit{Proc. IEEE SISY}, 2021.
\end{thebibliography}

\end{document}
